% ============================================================
% Packages
% ============================================================

\usepackage[
    a4paper,
    landscape,
    margin=10mm
]{geometry}

\usepackage{lmodern}
\usepackage{multicol}
\usepackage{xcolor}
\usepackage{listings}
\usepackage{fancyhdr}
\usepackage{graphicx}
\usepackage{xstring}

% ============================================================
% Layout
% ============================================================

\setlength{\columnsep}{3mm}
\setlength{\parindent}{0pt}
\setlength{\parskip}{0pt}

% ============================================================
% Colors
% ============================================================

\definecolor{codewhite}{RGB}{255,255,255}
\definecolor{codegray}{RGB}{245,245,245}
\definecolor{commentgray}{RGB}{100,100,100}
\definecolor{keywordblue}{RGB}{0,60,140}
\definecolor{stringred}{RGB}{140,30,30}

% ============================================================
% Code font
% ============================================================


\newcommand{\codefont}{%
    \ttfamily
    \fontsize{8pt}{9pt}\selectfont
}

% ============================================================
% Code listings
% ============================================================

\lstdefinestyle{python}{
    language=Python,
    basicstyle=\codefont,
    keywordstyle=\color{keywordblue}\bfseries,
    commentstyle=\color{commentgray},
    stringstyle=\color{stringred},
    backgroundcolor=\color{codewhite},
    numbers=none,
    frame=single,
    framerule=0.2pt,
    rulecolor=\color{black!25},
    breaklines=true,
    breakatwhitespace=true,
    columns=fullflexible,
    breakindent=0pt,
    keepspaces=true,
    showstringspaces=false,
    % Tabs are displayed as 2 spaces.
    tabsize=2,
    % Four spaces in the source are displayed as two spaces.
    %literate={\ \ }{{\ }}2,
    literate={\ \ }{{\ }}2 {~}{{\textasciitilde}}1,
    % Mark lines that have been wrapped around
    prebreak=\mbox{\textcolor{red}\textbackslash},
    aboveskip=2pt,
    belowskip=4pt
}

\lstdefinestyle{cpp}{
    language=C++,
    basicstyle=\codefont,
    keywordstyle=\color{keywordblue}\bfseries,
    commentstyle=\color{commentgray},
    stringstyle=\color{stringred},
    backgroundcolor=\color{codewhite},
    numbers=none,
    frame=single,
    framerule=0.2pt,
    rulecolor=\color{black!25},
    breaklines=true,
    breakatwhitespace=true,
    columns=fullflexible,
    breakindent=0pt,
    keepspaces=true,
    showstringspaces=false,
    % Tabs are displayed as 2 spaces.
    tabsize=2,
    % Four spaces in the source are displayed as two spaces.
    %literate={\ \ }{{\ }}2,
    literate={\ \ }{{\ }}2 {~}{{\textasciitilde}}1,
    % Mark lines that have been wrapped around
    prebreak=\mbox{\textcolor{red}\textbackslash},
    aboveskip=2pt,
    belowskip=4pt
}

\lstdefinelanguage{vim}{
    % list of keywords
    morekeywords={
        set, let
        map, nmap,
        filetype,
        on, off,
        autocmd,
        Plugin,
        call,
    },
    morecomment=[l]{"}, % l is for line comment
    morestring=[b]' % defines that strings are enclosed in double quotes
}

\lstdefinestyle{vimrc}{
    language=vim,
    basicstyle=\codefont,
    keywordstyle=\color{keywordblue}\bfseries,
    commentstyle=\color{commentgray},
    stringstyle=\color{stringred},
    backgroundcolor=\color{codewhite},
    numbers=none,
    frame=single,
    framerule=0.2pt,
    rulecolor=\color{black!25},
    breaklines=true,
    breakatwhitespace=true,
    columns=fullflexible,
    breakindent=0pt,
    keepspaces=true,
    showstringspaces=false,
    % Tabs are displayed as 2 spaces.
    tabsize=2,
    % Four spaces in the source are displayed as two spaces.
    %literate={\ \ }{{\ }}2,
    literate={\ \ }{{\ }}2 {~}{{\textasciitilde}}1,
    % Mark lines that have been wrapped around
    prebreak=\mbox{\textcolor{red}\textbackslash},
    aboveskip=2pt,
    belowskip=4pt
}

\lstdefinestyle{markdown}{
    language=html,
    basicstyle=\codefont,
    keywordstyle=\color{keywordblue}\bfseries,
    commentstyle=\color{commentgray},
    stringstyle=\color{stringred},
    backgroundcolor=\color{codewhite},
    numbers=none,
    frame=single,
    framerule=0.2pt,
    rulecolor=\color{black!25},
    breaklines=true,
    breakatwhitespace=true,
    columns=fullflexible,
    breakindent=0pt,
    keepspaces=true,
    showstringspaces=false,
    % Tabs are displayed as 2 spaces.
    tabsize=2,
    % Four spaces in the source are displayed as two spaces.
    %literate={\ \ }{{\ }}2,
    literate={\ \ }{{\ }}2 {~}{{\textasciitilde}}1,
    % Mark lines that have been wrapped around
    prebreak=\mbox{\textcolor{red}\textbackslash},
    aboveskip=2pt,
    belowskip=4pt
}

% ============================================================
% File heading
% ============================================================

\newcommand{\codefile}[1]{%
    \vspace{3pt}%
    \parbox{\linewidth}{%
        \raggedright
        {\small\bfseries\ttfamily\detokenize{#1}}%
    }%
    \par\vspace{1pt}%
}

% ============================================================
% Compact Table of Contents
% ============================================================

% We use LaTeX's normal .toc mechanism, but make the entries
% extremely compact.

\makeatletter

% Remove the normal "Contents" heading formatting.
% We don't want a "Contents" heading at all.
\renewcommand{\tableofcontents}{%
    \@starttoc{toc}%
}

% Make each TOC entry tiny and compact.
\renewcommand{\l@section}[2]{%
    \noindent
    \makebox[0.70\columnwidth][l]{%
        \ttfamily\scriptsize
        #1%
        \leaders\hbox to 1.0em{\hss.\hss}\hfill
    }%
    \ttfamily\scriptsize #2%
    \par
}

\makeatother

% ============================================================
% Include code file
%
% Every file automatically adds itself to the TOC.
% ============================================================

\newcommand{\mylisting}[2]{%
    \IfEndWith{#1}{.py}{%
        \lstinputlisting[style=python,inputpath=#2]{#1}%
    }{%
        \IfEndWith{#1}{.cpp}{%
            \lstinputlisting[style=cpp,inputpath=#2]{#1}%
        }{%
            \IfEndWith{#1}{.h}{%
                \lstinputlisting[style=cpp,inputpath=#2]{#1}%
            }{%
                \IfEndWith{#1}{.vimrc}{%
                    \lstinputlisting[style=vimrc,inputpath=#2]{#1}%
                }{%
                    \IfEndWith{#1}{.md}{%
                        \lstinputlisting[style=markdown,inputpath=#2]{#1}%
                    }{%
                        \PackageError{myfile}{Unknown file extension}{%
                            No listings style is defined for `#1'.%
                        }%
                    }%
                }%
            }%
        }%
    }%
}

\newcommand{\pyrivalfile}[1]{%
    \addcontentsline{toc}{section}{\protect\detokenize{#1}}%
    \codefile{#1}%
    \mylisting{#1}{\pyrivalroot}%
}

\newcommand{\customfile}[1]{%
    \addcontentsline{toc}{section}{\protect\detokenize{#1}}%
    \codefile{#1}%
    \mylisting{#1}{\customroot}%
}

% ============================================================
% Page style
% ============================================================

\pagestyle{fancy}
\fancyhf{}

\renewcommand{\headrulewidth}{0pt}

% Page number in the top-right corner.
% Lowered slightly so it is not cut off.
\fancyhead[R]{\raisebox{-10mm}{\textbf{\scriptsize\thepage}}}

% University name in the top-left corner.
% Lowered slightly so it is not cut off.
\fancyhead[L]{\raisebox{-10mm}{\textbf{\univabbr}}}
